\section{Flatness of the original normalized simultaneous model}

Let $k$ be a field, let $X$ be the original projective variety, and let
$I,J$ be ideal sheaves on its actual product with the polynomial line.
Assume that their model ideal logarithms are finite at every original
variety Gauss point. Write $B$ for the actual simultaneous blowup and
$\pi:N\to B$ for its actual generic-field normalization, with the original
composite structure map $g:N\to\operatorname{Spec}k[t]$.

\begin{lemma}[Surjectivity of the original normalized-model projection]
\label{local:original-normalized-pl-projection-surjective}
\uses{def:actual-original-normalized-simultaneous-model,thm:actual-original-pl-simultaneous-blowup-integral,scheme:actual-integral-normalization-surjective}
\lean{EnvelopeProperty.ProjSubvariety.normalizedSimultaneousProjection_surjective}
\leanok
The same actual projection $\pi:N\to B$ is surjective over any field $k$.
\end{lemma}
\begin{proof}
At the original trivial analytic point of $X$, the two assumed finite
logarithms prove integrality of the same simultaneous blowup $B$.
Apply the surjectivity theorem for the actual generic-field normalization
of this integral scheme. More explicitly, every point of $B$ lies in a
nonempty actual affine open with domain of sections $A$; lying-over for
the injective integral inclusion of $A$ into its integral closure in the
original function field gives a point of that same normalization above it.
The actual normalization chart triangle identifies its image with the
original point. These are precisely $N$ and $\pi$ in the stated construction.
\end{proof}

\begin{theorem}[Dominance of the original normalized model over any field]
\label{local:original-normalized-pl-model-dominant-any-field}
\uses{local:original-normalized-pl-projection-surjective,lem:pl-blowup-gauss-lift,berkovich:actual-gauss-point-dominance,def:actual-original-normalized-simultaneous-model}
\lean{EnvelopeProperty.ProjSubvariety.normalizedSimultaneousStructure_isDominant_anyField}
\leanok
The original map $g:N\to\operatorname{Spec}k[t]$ is dominant over any field
$k$. No projective normality assumption on $X$ is required.
\end{theorem}
\begin{proof}
Use the original trivial analytic point of $X$. The finite ideal logarithms
lift its actual positive-radius Gauss point through the original blowup's
off-centre isomorphism. Its whole prescribed polynomial-base point remains
the original Gauss point. That valuation has zero support: a nonzero
polynomial has a nonzero coefficient, whose contribution to its Gauss
value is strictly positive because the radius is positive. Thus the actual
support of the lifted point maps to the generic point of
$\operatorname{Spec}k[t]$. Its closure is the whole base, so the original
structure map of $B$ is dominant. The preceding surjectivity makes
$\pi$ dominant as well. Composition of these actual dominant maps proves
dominance of their same original composite $g$.
\end{proof}

\begin{theorem}[Flatness of the original normalized model over any field]
\label{local:original-normalized-pl-model-flat-any-field}
\uses{local:original-normalized-pl-model-dominant-any-field,thm:actual-original-normalized-simultaneous-model-normal,thm:proper-flattening-dedekind-native-flat}
\lean{EnvelopeProperty.ProjSubvariety.normalizedSimultaneousStructure_isFlat_anyField}
\leanok
The same actual original map $g:N\to\operatorname{Spec}k[t]$ is flat over
any field $k$. Its only model hypotheses are the stated finite original
Gauss ideal logarithms.
\end{theorem}
\begin{proof}
The constructed normalization $N$ is integral, the preceding theorem
proves actual dominance of $g$, and $k[t]$ is a Dedekind domain. Apply
the dominance-to-flatness theorem over a Dedekind domain to this very
map. Its coefficient argument applies on each nonempty original affine
open $U$ of $N$: its inclusion is dominant because $N$ is irreducible,
so the actual ring map $k[t]\to\Gamma(N,U)$ has kernel contained in the
zero nilradical of $k[t]$. This map is injective. Since $\Gamma(N,U)$ is
a domain, its actual $k[t]$-module is torsion free, hence flat over the
Dedekind domain $k[t]$. The affine flatness criterion and source locality
give flatness of the original $g$. Thus no general assertion that
normalization preserves flatness, and no replacement of the original
model or coefficient action, enters the argument.
\end{proof}
